\chapter*{\center \Large RECOMENDACIONES}
\addcontentsline{toc}{section}{\bfseries RECOMENDACIONES}
\markboth{RECOMENDACIONES}{RECOMENDACIONES}

Las recomendaciones que siguen se desprenden de los resultados del plan
de pruebas (Sección~\ref{sec:plan2}) y de la fase exploratoria que lo
precedió. Las primeras son metodológicas: valen para cualquier
continuación del trabajo, con estos escenarios o con otros. Las
segundas ordenan el trabajo siguiente.

\subsection*{Recomendaciones metodológicas}

\noindent\textbf{Barrer siempre la referencia fija en reparto, ciclo y
desfase, y ampliar la malla hasta un óptimo interior.} El primer
barrido de esta tesis variaba un único verde común para todas las
aproximaciones y dio por tuneado un programa que, en el corredor, era
10~s peor que el mejor reparto asimétrico; contra él, la mejora de IPPO
parecía de un $38$\,\% y era de un $6$\,\%. En el plan de pruebas el
barrido se amplió en cada caso hasta que el óptimo quedó en el interior
de la malla, con el verde mínimo de 10~s como única frontera aceptada.
Esa ampliación confirmó los programas anteriores del corredor (20~s de
verde de avenida, 10~s de transversal y 15~s de desfase), del corredor
con demanda doble y de la red, y fijó los de la red con demanda alta y
la malla. Conviene repetirla en todo escenario nuevo, con semillas
distintas de las de evaluación, antes de comparar nada con el fijo.
% fuente: plan2/barrido_<esc>.csv (fila de menor timeloss_prom y ejes barridos)

\noindent\textbf{Mantener las dos referencias, fija y actuada.} Un
resultado que solo supere al tiempo fijo no demuestra utilidad
práctica, porque el actuado está disponible de fábrica y no requiere
entrenamiento. El plan de pruebas mostró además que ninguna de las dos
basta sola: el actuado supera al fijo con demanda doble en el corredor
($-19.8$\,\%) y pierde contra él en la malla ($+21.8$\,\%). Ambas deben
tener las mismas fases y los mismos límites de verde que el agente.
% fuente: plan2/veredicto.csv (actuado)

\noindent\textbf{Juzgar la convergencia alrededor de la política
publicada, no solo al final de la serie.} El criterio de convergencia
del plan mira el último tercio de la serie de validación, pero la
política publicada puede estar antes. En la red, las dos bases de IPPO
\emph{vecinos} que funcionan publicaron los episodios 40 y 50 de 200,
cada uno un mínimo aislado de su serie (19.69~s frente a 21.29~s o más
en el resto de las validaciones), y la tabla las marca como
convergidas porque el último tercio es plano en un valor peor. Se
recomienda reportar, junto al veredicto, la distancia entre la
validación elegida y la mediana de la serie, y exigir que el mínimo se
repita en dos validaciones seguidas o usar más semillas de validación.
Para la única comparación que la tesis quiera afirmar como principal
(\emph{vecinos} frente a \emph{local}) convienen cinco semillas base en
vez de tres: en la red, una de las tres bases de \emph{vecinos} empata
con el fijo mientras las otras dos lo superan en más del 7\,\%. La
convergencia es además lo que deja abiertas cinco de las seis celdas
pendientes del plan: Q-learning en P4, Q-learning con la recompensa de
colas en P2, DQN en P5 y MAPPO en P3 y P4. Antes de dar por definitivo
su veredicto provisional conviene prolongar esos entrenamientos o sumar
bases, sin cambiar el resto del método.
% fuente: plan2/politica_red_vecinos_s*.json (val_serie, ep_elegido); plan2/veredicto.csv
% (convergio, convergen, veredicto); plan2/evaluacion_red_resumen.csv (ippo_vecinos_s2042, _s42, _s4042)

\noindent\textbf{Comprobar la semántica de la observación y de las fases
con pruebas automáticas.} Dos errores de esta clase aparecieron durante
el plan. El prototipo de la malla del 18.09.2026 tenía cruzadas las
fases este--oeste y norte--sur; la malla evaluada se construyó
corregida, y el barrido detecta ahora la fase de la avenida a partir de
la red. En la observación de IPPO, la ranura llamada
\texttt{cola\_mi\_salida} mide en realidad la cola del enlace que entra
desde el vecino; viene de la versión del 06.09.2026 y en la malla, que
es simétrica, da el mismo valor. Cambiar esa definición cambiaría todas
las políticas entrenadas, así que en esta tesis se documenta y no se
corrige; antes de cualquier entrenamiento nuevo conviene renombrarla o
corregirla y agregar una prueba que compare cada ranura con una lectura
directa de TraCI.

\noindent\textbf{Calibrar la demanda de estrés mirando todos los
cuellos de botella, no solo los semáforos.} El criterio heredado del
corredor con demanda doble buscaba el factor que llena el enlace entre
cruces. En la red, ese factor habría sido 2.5, pero la rotonda sin
semáforo se satura mucho antes: desde 1.3 veces la demanda nominal sus
ramales de un carril se llenan y cada vehículo que entra por ellos
pierde entre 400 y 700~s que ningún semáforo puede evitar. El caso P4
usa por eso 1.2 veces la demanda nominal. Todo escenario de estrés debe
comprobar antes que la congestión que se mide la pueden resolver los
controles que se comparan.
% fuente: tabla por factor en la cabecera de plan2/red_alta.rou.xml

\subsection*{Trabajo siguiente}

\noindent\textbf{Cerrar las celdas pendientes y ajustar DQN y MAPPO por
caso.} Al congelar los resultados, P1 a P5 tenían los ocho aprendices
evaluados salvo la ablación con la recompensa de colas en P5, que no se
corrió; es la única celda sin datos. Las otras cinco pendientes son de
convergencia (recomendación anterior). DQN y MAPPO usan los
hiperparámetros fijados en su prueba de humo en el cruce aislado. Un
ajuste corto por caso, con las semillas de validación y no con las de
evaluación, diría si la distancia de DQN a \emph{vecinos} (de 0.7 a
6.0~s en cuatro de los cinco casos) es del algoritmo o de los
hiperparámetros, y si MAPPO converge con más regularidad en las redes.
La variante \emph{cambio}, que penaliza cada cambio de fase, quedó
formulada y sin correr.
% fuente: plan2/veredicto.csv (ql_colas P5 sin datos; convergio); plan2/evaluacion_<esc>_resumen.csv
% (timeloss_prom de dqn_s4042 e ippo_vecinos publicada)

\noindent\textbf{Priorizar P6 sobre P7.} En la malla de nueve cruces
con demanda nominal ningún modelo supera al fijo coordinado en ajedrez; los
que reciben mensajes lo igualan (conclusión~5). En el corredor, en
cambio, la ganancia grande apareció cuando la demanda superó la
capacidad y el fijo dejó de bastar (P2). P6, la misma malla con
demanda alta, es el caso que puede decir si la coordinación aprendida
supera a un fijo coordinado en una malla; P7 agrega tamaño y fases
antes de haber contestado esa pregunta.

\noindent\textbf{Presupuestar P6 y P7 con el costo medido.} Los tiempos
se midieron en un AMD Ryzen 7 3700X de 8 núcleos y 16 hilos con 16~GB
de memoria, sin GPU, con un hilo por entrenamiento y varios
entrenamientos a la vez. La mediana del tiempo de reloj por episodio,
según el modelo, fue de 5.1 a 8.5~s en P1, de 12.7 a 16.6~s en P2, de
8.9 a 11.1~s en P3, de 9.8 a 12.1~s en P4 y de 23.5 a 30.7~s en P5;
DQN es el más caro en P1, P3 y P5. Un entrenamiento de 200
episodios, con sus validaciones, tardó entre 21 y 63 minutos por
semilla base en P1 a P4 y entre 95 y 115 minutos en P5. Si la demanda
alta encarece el episodio de la malla como en el corredor, donde
duplicarla lo multiplicó por 2.3, un entrenamiento de P6 tardaría entre
3.6 y 4.4~h por semilla base, del orden de 11 a 13~h de hilo por
modelo con tres bases. P7 (dieciséis cruces con tres carriles y giro
protegido) no se ha construido y su costo no se ha medido; al pasar de
tres a nueve semáforos el episodio se encareció 2.7 veces, y con cuatro
fases el tabular choca además con el techo de estados del cruce con
giros protegidos. Estas cifras de P6 y P7 son estimaciones a partir de
los tiempos medidos, no mediciones.
% fuente: columna segundos de plan2/resultados_<esc>_<alg>_s<base>.csv (mediana por corrida; se
% excluyen las corridas de confirmacion de ippo_spill en curso); duracion por trabajo: "tiempos" en
% plan2/cola/estado.json (train_* de P1-P4: 1268 a 3788 s; train_malla3_*: 5704 a 6917 s);
% P6 = 2.3 x duracion de P5 (estimacion)

\noindent\textbf{Probar el término de \textit{spillback} donde puede
actuar.} En las dos redes y en la malla el término nunca se activa, y
en el corredor con demanda doble no mejora a \emph{vecinos}.
Queda sin probar el rango intermedio: una demanda media en el corredor
(entre 1.3 y 1.5 veces la nominal, donde la calibración encontró el
enlace lleno en cero y en una de tres semillas) o un umbral más bajo
que la mitad de las plazas en la red. Solo ahí se puede saber si el
término aporta algo que la mensajería no aporte.

\noindent\textbf{Entrenar la coordinación con un canal imperfecto.} Las
políticas de U4 se entrenaron con mensajes instantáneos y, al
evaluarlas con latencia o pérdida en la fase exploratoria, usaban el
mensaje viejo como si fuera actual: el retraso global apenas cambiaba
pero la progresión se deshacía. La edad del mensaje ya está en la
observación; basta entrenar con latencia y pérdida aleatorias por
episodio para que el agente aprenda a descontarla, y evaluar con el
protocolo del plan. Es el paso previo a cualquier despliegue sobre una
red de comunicación real.

\noindent\textbf{Llevar el protocolo al cruce real de Lima.} La
importación desde OpenStreetMap del cruce de la Av.\ Paseo de la
República con la Av.\ 28 de Julio se hizo el 01.12.2025 con
\texttt{osmWebWizard}, pero la red no se ha limpiado en
\texttt{netedit}, su demanda es la aleatoria del asistente y no ha
pasado por este protocolo. Conviene completarla con el mismo método de
los casos P1 a P5: barrido del fijo, actuado con las mismas fases,
tres semillas base y evaluación sobre las mismas semillas. La demanda
del caso real debe fijarse cerca de la capacidad y no muy por encima,
porque con demanda muy superior ningún control mueve el
\textit{throughput}: todos sirven el máximo que admite el carril.

\noindent\textbf{Evaluar con llegadas aleatorias.} Los escenarios
actuales insertan los vehículos de forma equiespaciada, de modo que la
única variabilidad entre semillas es el comportamiento de los
conductores. Repetir la evaluación con llegadas de Poisson daría una
medida más realista de la robustez de cada control, a costa de rehacer
el ajuste de los programas de referencia.
